\documentclass[12pt]{article}
\usepackage[utf8]{inputenc}
\usepackage{amsmath}
\usepackage{graphicx}
\usepackage{hyperref}
\usepackage{listings}
\usepackage{float}

\setlength{\parindent}{0pt}
\title{ERTS: Assignment 2}
\author{
Simon Møller, student number: 202204535 \\ 
Mark Myrtoft Bøgelund, student number: 202206634
}
\date{Date: \today}
\lstset{
    language=C++,
    basicstyle=\ttfamily\footnotesize,
    frame=single,
    breaklines=true
}

\newcommand{\codeplaceholder}[1]{%
\begin{figure}[h]
    \centering
    \fbox{%
        \begin{minipage}[c][30mm][c]{0.82\textwidth}
            \centering
            \texttt{[Code snippet placeholder]}\\[2mm]
            \textit{#1}
        \end{minipage}%
    }
    \caption{#1}
\end{figure}%
}

\newcommand{\resultplaceholder}[1]{%
\begin{figure}[h]
    \centering
    \fbox{%
        \begin{minipage}[c][45mm][c]{0.82\textwidth}
            \centering
            \textit{Placeholder for #1.}
        \end{minipage}%
    }
    \caption{#1}
\end{figure}%
}

\begin{document}
\maketitle

\section*{Introduction}

The purpose of this assignment is to learn the design flow for a system on a
programmable chip. Using Vivado, SDK, and HLS, we explore the relationship
between the Zynq Processing System (PS) and Programmable Logic (PL), interface
drivers, command-line control, custom IP, and hardware acceleration.

We began by following the setup instructions to configure the ZYBO board,
connect it to the development computer, and prepare the required Xilinx design
tools. This was initially quite challenging, since our pc runs on Windows 11,
and the old software was not fully compatible. Eventually, we got it working.
This provided the hardware and software platform for the exercises that follow.

\newpage
\section*{Exercise 3: Create a Console Application}

In this exercise, we create a small console application for the ZYBO board. It
receives commands over USB-UART and uses them to control the LEDs through the
switches and timer.

\subsection*{Command interpreter}

\codeplaceholder{Console command interpreter}

The application will read a command from the terminal and use a \texttt{switch}
statement to select the requested operation.

\subsection*{Command 1: Switches to LEDs}

\codeplaceholder{Reading the switches and writing the LED output}

\begin{figure}[H]
    \centering
    \includegraphics[width=0.5\textwidth]{\detokenize{../images/E3 - Switch and LED.jpg}}
    \caption{The LEDs reflect the binary state of the switches.}
\end{figure}

With command \texttt{1}, the LEDs light up according to the switch positions.

\subsection*{Command 2: Timer-driven binary counter}

\codeplaceholder{XScuTimer initialization and one-second counter}

\begin{figure}[H]
    \centering
    \includegraphics[width=0.5\textwidth]{\detokenize{../images/E3 - Counter and LED 1.jpg}}
    \caption{The terminal displays the counter while the LEDs show its binary value.}
\end{figure}

\begin{figure}[H]
    \centering
    \includegraphics[width=0.5\textwidth]{\detokenize{../images/E3 - Counter and LED 2.jpg}}
    \caption{The counter continues increasing and is reflected on the LEDs.}
\end{figure}

With command \texttt{2}, the timer increases the counter once per second. The
LEDs display the counter as a four-bit binary value.

\subsection*{Discussion}

This exercise gives us a first working connection between software running on
the processing system and hardware peripherals on the ZYBO board. It also
prepares the structure we will use in the later exercises.

\newpage
\section*{Exercise 4: Create a Matrix Multiplication}

In this exercise, we implement a software matrix multiplication on the ZYBO
board. The program uses two 4x4 matrices, calculates the product, and reports
the execution time through the USB-UART interface.

\subsection*{Matrix multiplication}

The first matrix contains the values from 1 to 16. For the second matrix, the
assignment gives $B^T$, so if:
\[
    B^T =
    \left(
    \begin{bmatrix}
            1 & 1 & 1 & 1 \\
            2 & 2 & 2 & 2 \\
            3 & 3 & 3 & 3 \\
            4 & 4 & 4 & 4
        \end{bmatrix}
    \right)^T
    =
    \begin{bmatrix}
        1 & 2 & 3 & 4 \\
        1 & 2 & 3 & 4 \\
        1 & 2 & 3 & 4 \\
        1 & 2 & 3 & 4
    \end{bmatrix}.
\]

Then:
\[
    B =
    \begin{bmatrix}
        1 & 1 & 1 & 1 \\
        2 & 2 & 2 & 2 \\
        3 & 3 & 3 & 3 \\
        4 & 4 & 4 & 4
    \end{bmatrix}
\]
We used the matrices $A$ and $B$ as defined above for the multiplication, which
is performed using nested loops, where each element in the result is calculated
as the sum of four products.

Following the operations shown in the figure below, which we took from the
assignment:

\begin{figure}[H]
    \centering
    \includegraphics[width=0.5\textwidth]{\detokenize{../images/E4 - Assignment operations.png}}
    \caption{Matrix multiplication operations for calculating the elements of $P$.}
\end{figure}

We made the multiplication and transposition in the same loops. This was
confusing at first, but we eventually got it and arrived at the correct result.

\begin{lstlisting}[language=C,caption={Matrix initialization and software multiplication}]
void setInputMatrices(VectorArray A, VectorArray B)
{...}
void multiMatrixSoft(VectorArray A, VectorArray B, VectorArray P)
{
    int Arow, Brow, col;
    unsigned int sum;
    for (Arow = 0; Arow < MSIZE; Arow++) {
        for (Brow = 0; Brow < MSIZE; Brow++) {
            sum = 0;
            for (col = 0; col < MSIZE; col++)
                sum += A[Arow].comp[col] * B[Brow].comp[col];
            P[Arow].comp[Brow] = sum;
        }
    }
}
\end{lstlisting}

The full implementation and setup can be found in the accompanying source code
files.

The result of the matrix multiplication is shown in the figure below.

\begin{figure}[H]
    \centering
    \includegraphics[width=0.5\textwidth]{\detokenize{../images/E4 - Matrix Mult.jpg}}
    \caption{Matrix multiplication result and execution time.}
\end{figure}

The result matrix is displayed correctly, and the software implementation takes
1102 clock ticks in this run. This result will later be used for comparison
with the hardware-accelerated version.

\newpage
\section*{Exercise 5: Create a Hardware IP Acceleration}

In this exercise, we use a custom hardware IP to accelerate the matrix
multiplication from Exercise 4. The IP performs four byte-wise multiplications
and adds the results before returning a single matrix element.

\subsection*{Matrix IP}

The matrix IP receives two packed 32-bit values through its registers. Each
value contains four matrix elements, so the IP can multiply the corresponding
bytes in parallel and return their sum. The result is read through a third
memory-mapped register.

\subsection*{Hardware matrix multiplication}

\begin{lstlisting}[language=C,caption={Writing matrix rows to the IP and reading each result}]
void multiMatrixHard(VectorArray A, VectorArray B, VectorArray P)
{
    for (int row = 0; row < MSIZE; row++) {
        for (int col = 0; col < MSIZE; col++) {
            MATRIX_IP_mWriteReg(XPAR_MATRIX_IP_0_S00_AXI_BASEADDR,
                MATRIX_IP_S00_AXI_SLV_REG0_OFFSET, A[row].vect);
            MATRIX_IP_mWriteReg(XPAR_MATRIX_IP_0_S00_AXI_BASEADDR,
                MATRIX_IP_S00_AXI_SLV_REG1_OFFSET, B[col].vect);
            P[row].comp[col] = MATRIX_IP_mReadReg(
                XPAR_MATRIX_IP_0_S00_AXI_BASEADDR,
                MATRIX_IP_S00_AXI_SLV_REG2_OFFSET);
        }
    }
}
\end{lstlisting}

The hardware implementation uses the same matrix structure as the software
version. Instead of calculating the four products in the processor, each row of
$A$ and each row of $B$ are written to the IP. The returned value is then
stored in the corresponding element of $P$.

\subsection*{Timing comparison}

We measured elapsed clock ticks around each software and hardware function call
using the processor timer:
\begin{lstlisting}[language=C,caption={Measuring software and hardware execution time}]
XTime_GetTime(&startTime);
multiMatrixSoft(aInst, bInst, pInst);
XTime_GetTime(&endTime);
softTime = endTime - startTime;

XTime_GetTime(&startTime);
multiMatrixHard(aInst, bInst, pInst);
XTime_GetTime(&endTime);
hardTime = endTime - startTime;
\end{lstlisting}

Running the software and hardware matrix multiplication functions produces the
following comparison:
\begin{figure}[H]
    \centering
    \includegraphics[width=0.5\textwidth]{\detokenize{../images/E5 - Matrix mult soft and hard.jpg}}
    \caption{Comparison of the software and hardware matrix multiplication results.}
\end{figure}
\newpage
Both implementations produce the same result. The software version took 1253
clock ticks, compared with 3981 ticks for the hardware version. The processing
system runs at 650 MHz, while the hardware IP in the programmable logic runs at
100 MHz. The hardware time also includes transfers to and from the IP. If the
hardware IP ran at 650 MHz, its ideal scaled time would be:
\[
    3981 \cdot \frac{100}{650} \approx 612\ \text{clock ticks}.
\]
Compared with the software result, this corresponds to an ideal speedup of
\[
    \frac{1253}{612} \approx 2.05\times.
\]
The practical speedup would be lower because the transfer overhead remains.

\newpage
\section*{Exercise 7: HLS exercise implementing a Neural Network}

In this exercise, we implement a configurable two-layer neural network in C++
and use Vivado HLS to synthesize it as an IP core. The network uses fixed-point
integer arithmetic, with weights scaled by $2^{15}$, and applies a ReLU
activation to the hidden layer.

\subsection*{Neural network implementation}

The network has 196 inputs, 32 hidden neurons, and 10 outputs. The weights are
represented in fixed-point form with 15 fractional bits. After each dot
product, the result is shifted back by 15 bits; ReLU is applied to the
hidden-layer values.

\begin{lstlisting}[caption={Fixed-point scaling and ReLU in the hidden layer}]
int value = static_cast<int>(sum >> BITS);
neurons[i] = value < 0 ? 0 : value;
\end{lstlisting}

The input and weight arrays are controlled by the processor through AXI4-Lite.
Switches and LEDs use direct ports, with the switch value passed through to the
LEDs.

\begin{lstlisting}[caption={HLS interfaces for processor data and board I/O}]
#pragma HLS INTERFACE s_axilite port=input bundle=CTRL
#pragma HLS INTERFACE s_axilite port=weightsI bundle=CTRL
#pragma HLS INTERFACE s_axilite port=weightsO bundle=CTRL
#pragma HLS INTERFACE s_axilite port=output bundle=CTRL
#pragma HLS INTERFACE ap_none port=inSwitch
#pragma HLS INTERFACE ap_none port=outLeds

*outLeds = inSwitch;
\end{lstlisting}

The provided `neuralNet.h` defines the function interface, and `testBench.cpp`
supplies test inputs and checks the C model before synthesis.

\subsection*{HLS results}

We compared the design with and without pipelining the inner loops. In the
pipelined version, an initiation interval of one is requested:

\begin{lstlisting}[caption={Pipeline directive used on the inner loops}]
#pragma HLS PIPELINE II=1
\end{lstlisting}

\begin{figure}[H]
    \centering
    \includegraphics[width=0.82\textwidth]{\detokenize{../images/E7 - Simulation Synthesis.jpg}}
    \caption{Vivado HLS simulation and synthesis results.}
\end{figure}

\begin{figure}[H]
    \centering
    \includegraphics[width=0.82\textwidth]{\detokenize{../images/E7 - Simulation Log.jpg}}
    \caption{Simulation log for the neural network testbench.}
\end{figure}

\begin{figure}[H]
    \centering
    \includegraphics[width=0.82\textwidth]{\detokenize{../images/E7 - Comparison Report.jpg}}
    \caption{Comparison of the pipelined and non-pipelined HLS implementations.}
\end{figure}

\subsection*{ZYBO integration and timing}

The synthesized IP was exported as VHDL version 1.0.1 and integrated into the
Vivado ZYBO design. The board photos document the switch-to-LED pass-through.
The SDK test compares software and hardware outputs and measures input, weight,
compute, and output-transfer times separately, making the data-transfer cost
visible.

\begin{figure}[H]
    \centering
    \begin{minipage}[t]{0.48\textwidth}
        \centering
        \includegraphics[width=\textwidth]{\detokenize{../images/E7 - Zybo board LEDS and switches 1.jpg}}
        \caption{ZYBO switch and LED verification.}
    \end{minipage}\hfill
    \begin{minipage}[t]{0.48\textwidth}
        \centering
        \includegraphics[width=\textwidth]{\detokenize{../images/E7 - Zybo board LEDS and switches 2.jpg}}
        \caption{A second switch and LED test state.}
    \end{minipage}
\end{figure}

The complete source code is provided separately with the assignment and is not
repeated here.

\end{document}